%%% variant: default
\section{Wärmelehre}

\subsection{Wärmekapazität}

Die Wärmekapazität $C$ gibt an, wie viel Wärme $Q$ für eine Temperaturänderung
$\Delta T$ nötig ist.

\[
Q = C \cdot \Delta T
\]

\subsection{Rechenbeispiel}

Ein Körper mit einer Wärmekapazität von $500\,\mathrm{J/K}$ wird um
$2\,\mathrm{K}$ erwärmt. Wie viel Wärme nimmt er auf?



%%% variant: instructor
\section{Wärmelehre}

\subsection{Wärmekapazität}

Die Wärmekapazität $C$ gibt an, wie viel Wärme $Q$ für eine Temperaturänderung
$\Delta T$ nötig ist.

\[
Q = C \cdot \Delta T
\]

\subsection{Rechenbeispiel}

Ein Körper mit einer Wärmekapazität von $500\,\mathrm{J/K}$ wird um
$2\,\mathrm{K}$ erwärmt. Wie viel Wärme nimmt er auf?


Lösung: $Q = 500\,\mathrm{J/K} \cdot 2\,\mathrm{K} = 1000\,\mathrm{J}$.



%%% variant: handout
\section{Wärmelehre}

\subsection{Wärmekapazität}

Die Wärmekapazität $C$ gibt an, wie viel Wärme $Q$ für eine Temperaturänderung
$\Delta T$ nötig ist.

\[
Q = C \cdot \Delta T
\]

\subsection{Rechenbeispiel}

Ein Körper mit einer Wärmekapazität von $500\,\mathrm{J/K}$ wird um
$2\,\mathrm{K}$ erwärmt. Wie viel Wärme nimmt er auf?



\par\noindent\textbf{Vortragsnotiz:} Vor der Auflösung nach der Einheit fragen.\par

%%% variant: slides
\begin{frame}[t]
\frametitle{Wärmelehre}

\end{frame}

\begin{frame}[t]
\frametitle{Wärmekapazität}

Die Wärmekapazität $C$ gibt an, wie viel Wärme $Q$ für eine Temperaturänderung
$\Delta T$ nötig ist.

\[
Q = C \cdot \Delta T
\]

\end{frame}

\begin{frame}[t]
\frametitle{Rechenbeispiel}
\note{Vor der Auflösung nach der Einheit fragen.}

Ein Körper mit einer Wärmekapazität von $500\,\mathrm{J/K}$ wird um
$2\,\mathrm{K}$ erwärmt. Wie viel Wärme nimmt er auf?



\end{frame}

%%% variant: slides-dozent-h1-frame
\begin{frame}[t]
\frametitle{Wärmelehre}

\end{frame}

\begin{frame}[t]
\frametitle{Wärmekapazität}

Die Wärmekapazität $C$ gibt an, wie viel Wärme $Q$ für eine Temperaturänderung
$\Delta T$ nötig ist.

\[
Q = C \cdot \Delta T
\]

\end{frame}

\begin{frame}[t]
\frametitle{Rechenbeispiel}
\note{Vor der Auflösung nach der Einheit fragen.}

Ein Körper mit einer Wärmekapazität von $500\,\mathrm{J/K}$ wird um
$2\,\mathrm{K}$ erwärmt. Wie viel Wärme nimmt er auf?


Lösung: $Q = 500\,\mathrm{J/K} \cdot 2\,\mathrm{K} = 1000\,\mathrm{J}$.



\end{frame}

%%% variant: slides-h1-hide-level-1
\begin{frame}[t]
\note{Vor der Auflösung nach der Einheit fragen.}

\par\medskip\noindent\textbf{Wärmekapazität}\par


Die Wärmekapazität $C$ gibt an, wie viel Wärme $Q$ für eine Temperaturänderung
$\Delta T$ nötig ist.

\[
Q = C \cdot \Delta T
\]

\par\medskip\noindent\textbf{Rechenbeispiel}\par


Ein Körper mit einer Wärmekapazität von $500\,\mathrm{J/K}$ wird um
$2\,\mathrm{K}$ erwärmt. Wie viel Wärme nimmt er auf?



\end{frame}

